%
% input_text_reconciliation.tex
%
% A Z specification of the lux input_text honour/commit reconciliation
% state machine: the buffer/editing/commit discipline that decides, on a
% Hub/Display round trip with a DELAYED echo, when the local buffer defers
% the Hub value and when it honours it.  Formalises the second
% reconciliation-latency defect in input_text (Copilot MEDIUM, ``Refocus
% loses commit echo'') after the first (pipelined-edits clobber) motivated
% the commit-on-idle redesign.
%
% This model is the shared discipline of the continuous-edit reconciliation
% arbiters.  It is type-agnostic -- its carrier [VALUE] stands for the text an
% input_text holds or the float a slider holds -- so it governs BOTH arbiters
% unchanged; the shared-abstraction extraction is deferred to color_picker.
%
% Source of record:
%   src/punt_lux/display/renderers/imgui/input_text_selection.py
%       -- InputTextArbiter: resolve (honour-or-defer), observe, commit, release
%   src/punt_lux/display/renderers/input_text_renderer.py
%       -- InputTextRenderer.render: the per-frame active/deactivate/fire
%          decision over an ImGui input_text widget
%   src/punt_lux/display/renderers/imgui/slider_selection.py
%       -- SliderArbiter: the float sibling of InputTextArbiter, same
%          resolve/observe/commit/release honour-or-defer control flow
%   src/punt_lux/display/renderers/slider_renderer.py
%       -- SliderRenderer.render: the per-frame active/deactivate/fire
%          decision over an ImGui slider_float / slider_int widget
%   src/punt_lux/scene/widget_state.py
%       -- WidgetState: the per-scene buffer/editing slots (input and slider
%          triples) and their removal clear (discard_for)
%
% Type-check:  fuzz -t docs/input_text_reconciliation.tex
% Model-check (see Section "Verification" for the exact goal predicates):
%   probcli docs/input_text_reconciliation.tex -model_check -p DEFAULT_SETSIZE 2
%   probcli docs/input_text_reconciliation.tex -model_check -nodead \
%       -goal "lost=ztrue" -p DEFAULT_SETSIZE 2
%   probcli docs/input_text_reconciliation.tex -model_check -nodead \
%       -goal "editing=ztrue & edited=zfalse" -p DEFAULT_SETSIZE 2
%   probcli docs/input_text_reconciliation.tex -model_check -nodead \
%       -goal "fires>1" -p DEFAULT_SETSIZE 2
%   probcli docs/input_text_reconciliation.tex -model_check -nodead \
%       -goal "clobbered=ztrue" -p DEFAULT_SETSIZE 2
%
% Formatting follows the fuzz house rules: \quad~ for continuation
% indent (never \t1..\t9), schema lines under 80 columns, predicates
% broken at \land/\lor/\implies with the operator starting the
% continuation line.
%

\documentclass[a4paper,10pt,fleqn]{article}
\usepackage[margin=1in]{geometry}
\usepackage{fuzz}
\usepackage[colorlinks=true,linkcolor=blue,citecolor=blue,urlcolor=blue]{hyperref}

\begin{document}

\title{The lux input\_text Reconciliation Discipline: a Z Specification}
\author{Formal model of the honour / defer / commit-echo state machine}
\date{July 2026}
\maketitle

% ============================================================
\section{Introduction}
% ============================================================

An \texttt{input\_text} carries one Hub-authoritative \texttt{value}, replicated
to the Display and rendered every frame by an ImGui text widget that keeps its
\emph{own} editable buffer.  The renderer must reconcile the two.  While the
user is not editing, the field \emph{honours} the Hub: each frame renders the
replicated value, so an agent-driven change appears on the next frame, in the
manner of a checkbox.  While the user \emph{is} editing, the local buffer is
authoritative and a Hub-driven value is \emph{deferred}, so that two edits in
flight before a round trip returns cannot clobber the text under the user's
fingers.  Exactly one \texttt{ValueChanged} fires when an edit commits --- on
blur or Enter --- never per keystroke.

The subtlety the design must survive is the \emph{echo-latency window}.  The
Hub and the Display are separate processes, potentially on separate machines.
When the field commits a value, that value travels to the Hub, the Hub installs
it, and only then does it come back round as the replicated \texttt{value} the
Display renders.  Until that echo arrives, the replicated value the Display
holds is the \emph{pre-echo} one: the value from before the commit.  This
document calls that returning value the \emph{echo}, and the interval between a
commit and its echo the \emph{echo-latency window}.

Two defects of the same class have surfaced across two review rounds, and that
recurrence is why the discipline is modelled here rather than reasoned about
informally.  The first --- pipelined edits clobbering live text --- motivated
the commit-on-idle redesign, whose author held that ``the race is gone by
construction''.  A residual race was then found immediately, so the informal
argument did not hold.  The second (the residual) is the one this document is
built around: \emph{a re-focus during the echo-latency window loses the
committed value}.

The residual runs as follows.  On commit, the renderer drops its buffer and the
field returns to honouring the Hub; but the Hub still holds the pre-echo value,
so an idle frame renders that pre-echo value, not the value just committed.  If
the user re-focuses the field during this window, the current renderer begins
deferring on the mere act of focus: it snapshots the \emph{displayed} text ---
which is the stale pre-echo value --- as the buffer and sets the editing flag.
The arriving echo of the just-committed value is now deferred for the whole
focus session, because the field is editing; and subsequent typing commits from
that stale base, overwriting the value that had just committed.  On separate
machines this is routine, not exotic.

The model isolates the two decisions that, taken together, produce the loss:
\emph{when} the field begins deferring (on focus, or only on a real edit), and
\emph{what} an idle field honours during the echo-latency window (the raw
pre-echo Hub value, or the value it optimistically knows it committed).  It
proves that the corrected design --- defer only after a real edit, and honour
the committed value locally until the Hub echo confirms it --- admits no losing
trace, and that the two tempting partial fixes each still do.

This is a finite state machine over one field and a small fixed set of text
values, so ProB can enumerate every interleaving of focus, keystroke, commit,
echo, and re-focus, and either confirm the invariants or produce the exact
offending trace.

% ============================================================
\section{Basic Types}
% ============================================================

\begin{zed}
  [VALUE]
\end{zed}

\texttt{VALUE} is the set of text values the field may hold.  Two values suffice
to exhibit the loss: an initial and pre-echo value, and a distinct committed
value whose echo has not yet arrived.  Three exercise the interleaving of a
second edit against an in-flight echo more thoroughly.  The carrier is bounded
by ProB's \texttt{DEFAULT\_SETSIZE}, so no explicit cardinality constant is
needed.

% ============================================================
\section{Free Types and Constants}
% ============================================================

\begin{zed}
  ZBOOL ::= ztrue | zfalse
\end{zed}

A two-valued flag written as a free type rather than the toolkit \texttt{bool},
so ProB animates it as a small enumerated set.

\begin{zed}
  MAXFIRE == 2
\end{zed}

The fire counter saturates at \texttt{MAXFIRE}.  One edit session may fire at
most once; a count of two is already a witness to a repeated-fire defect, so
counting no higher keeps the state space small without hiding a violation.

\begin{axdef}
  v0 : VALUE
\end{axdef}

\texttt{v0} is the value the field starts on --- both the initial displayed text
and the initial Hub value.  A committed value distinct from \texttt{v0} is what
the echo-latency window is about: after committing it, the Hub still holds
\texttt{v0} until the echo lands.

% ============================================================
\section{State}
% ============================================================

The state is flat: one schema, eleven components.  Its predicate carries only
\emph{structural} coherence --- that a field can defer or be mid-edit only while
it is focused, and that the fire counter stays in range.  These hold in every
design, correct or buggy, so they act as coherence constraints, not as the
safety properties under test.  The safety properties (Section~\ref{sec:inv}) are
deliberately \emph{absent} here, so that a violating state remains reachable and
ProB can find it.

\begin{schema}{Field}
  disp : VALUE \\
  hub : VALUE \\
  focused : ZBOOL \\
  editing : ZBOOL \\
  edited : ZBOOL \\
  pending : ZBOOL \\
  committed : VALUE \\
  fires : 0 \upto MAXFIRE \\
  lost : ZBOOL \\
  clobbered : ZBOOL
\where
  editing = ztrue \implies focused = ztrue \\
  edited = ztrue \implies focused = ztrue
\end{schema}

\texttt{disp} is the text ImGui currently renders --- the value a commit would
fire.  \texttt{hub} is the Hub-authoritative replicated value.  \texttt{focused}
records whether the widget is being interacted with.  \texttt{editing} is the
\emph{defer} flag: while it is set the buffer is authoritative and the Hub value
is ignored.  \texttt{edited} is the \emph{real-edit} witness: whether a genuine
keystroke has changed the text since the current focus began; it distinguishes a
field the user is merely focused on from one being edited, and it is the
component the honour property is stated over.

\texttt{pending} is the echo-latency flag: set from a commit until the echo
lands, it marks the window in which the Hub has not yet caught up to the
committed value.  \texttt{committed} is the value most recently committed --- the
value the field optimistically displays while \texttt{pending}, and the value
whose durability is under test.  \texttt{fires} counts commit fires since the
current focus began, the window in which at most one fire is legitimate.

The last two components are pure bookkeeping for the invariant checks, not part
of the renderer's own state.  \texttt{lost} latches \texttt{ztrue} the moment an
edit session begins from a stale pre-echo base --- the durability failure.
\texttt{clobbered} latches \texttt{ztrue} the moment an echo overwrites live
post-edit text --- the no-clobber failure.

% ============================================================
\section{Initialization}
% ============================================================

\begin{schema}{Init}
  Field'
\where
  disp' = v0 \\
  hub' = v0 \\
  focused' = zfalse \\
  editing' = zfalse \\
  edited' = zfalse \\
  pending' = zfalse \\
  committed' = v0 \\
  fires' = 0 \\
  lost' = zfalse \\
  clobbered' = zfalse
\end{schema}

A single initial state: the field unfocused and not editing, display and Hub
both on \texttt{v0}, no commit outstanding, nothing fired, neither failure
latched.

% ============================================================
\section{The Focus Session}
% ============================================================

A focus session runs from the user focusing the field to the field losing focus.
Focusing begins a fresh session: it clears the real-edit witness and the fire
count, and --- this is the first half of the corrected design --- it does
\emph{not} begin deferring.  A focused field that has not yet been edited still
honours the Hub (Section~\ref{sec:frames}), so an echo arriving mid-focus can
still reach the display.  This is what distinguishes the corrected design from
the current one, whose focus begins deferring at once
(Section~\ref{sec:fidelity}, defect~(i)).

\begin{schema}{Focus}
  \Delta Field
\where
  focused = zfalse \\
  focused' = ztrue \\
  editing' = zfalse \\
  edited' = zfalse \\
  fires' = 0 \\
  disp' = disp \\
  hub' = hub \\
  pending' = pending \\
  committed' = committed \\
  lost' = lost \\
  clobbered' = clobbered
\end{schema}

A keystroke changes the text.  It is the second half of the corrected design:
deferring begins here, on the \emph{first} real edit of the session, not on
focus.  The new text \texttt{e?} models whatever the user typed.  The durability
latch reads the state at the instant deferring begins --- the transition
\texttt{editing} \texttt{zfalse} $\to$ \texttt{ztrue} --- and fires if that base
is stale: a commit is still in flight (\texttt{pending}), yet the base the edit
builds on is not the committed value.  Under the corrected honour rule this can
never happen (Section~\ref{sec:inv}); the latch is here so that the fidelity
variants, which weaken that rule, have somewhere to trip it.

\begin{schema}{Keystroke}
  \Delta Field \\
  e? : VALUE
\where
  focused = ztrue \\
  lost' = \\
  \quad~ \IF editing = zfalse \land pending = ztrue \land disp \neq committed \\
  \quad~ \THEN ztrue \ELSE lost \\
  editing' = ztrue \\
  edited' = ztrue \\
  disp' = e? \\
  focused' = focused \\
  hub' = hub \\
  pending' = pending \\
  committed' = committed \\
  fires' = fires \\
  clobbered' = clobbered
\end{schema}

A commit is the deactivation of a field that was being edited --- blur or Enter
after a real edit.  It fires exactly once, records the committed value, and
opens the echo-latency window (\texttt{pending}).  The Hub is \emph{not} updated
here: the echo is delayed, and arrives only later, in \texttt{HubEcho}.  The
display keeps the committed value: because \texttt{pending} is now set, the
following honour frame renders \texttt{committed}, not the stale Hub value.

\begin{schema}{Commit}
  \Delta Field
\where
  focused = ztrue \\
  editing = ztrue \\
  committed' = disp \\
  pending' = ztrue \\
  fires' = \IF fires < MAXFIRE \THEN fires + 1 \ELSE MAXFIRE \\
  focused' = zfalse \\
  editing' = zfalse \\
  edited' = zfalse \\
  disp' = disp \\
  hub' = hub \\
  lost' = lost \\
  clobbered' = clobbered
\end{schema}

A field may also lose focus without ever having been edited: the user focused it
and moved away.  This deactivation fires nothing --- there is no edit to commit
--- and is the reason a fire count is per session rather than global.

\begin{schema}{BlurNoEdit}
  \Delta Field
\where
  focused = ztrue \\
  editing = zfalse \\
  focused' = zfalse \\
  edited' = zfalse \\
  editing' = zfalse \\
  disp' = disp \\
  hub' = hub \\
  pending' = pending \\
  committed' = committed \\
  fires' = fires \\
  lost' = lost \\
  clobbered' = clobbered
\end{schema}

% ============================================================
\section{The Hub Value}
% ============================================================

Two operations move the Hub value.  \texttt{HubEcho} is the delayed echo of a
commit: the committed value finally arrives at the Hub and the echo-latency
window closes.  From here on an honour frame renders \texttt{hub}, which now
equals \texttt{committed}, so the optimistic display and the real Hub value have
converged.

\begin{schema}{HubEcho}
  \Delta Field
\where
  pending = ztrue \\
  hub' = committed \\
  pending' = zfalse \\
  disp' = disp \\
  focused' = focused \\
  editing' = editing \\
  edited' = edited \\
  committed' = committed \\
  fires' = fires \\
  lost' = lost \\
  clobbered' = clobbered
\end{schema}

\texttt{AgentPush} is an agent-driven change to the Hub value while the field is
settled --- idle, not editing, no commit outstanding.  It is the checkbox-style
honour case: a value pushed from outside must appear on the next honour frame.
It is guarded to the settled state so that it exercises honour without
entangling with the echo-latency window.

\begin{schema}{AgentPush}
  \Delta Field \\
  w? : VALUE
\where
  focused = zfalse \\
  editing = zfalse \\
  pending = zfalse \\
  hub' = w? \\
  disp' = disp \\
  focused' = focused \\
  editing' = editing \\
  edited' = edited \\
  pending' = pending \\
  committed' = committed \\
  fires' = fires \\
  lost' = lost \\
  clobbered' = clobbered
\end{schema}

% ============================================================
\section{The Render Frames}
\label{sec:frames}
% ============================================================

Each frame in which the field is \emph{not} deferring honours: it drives the
display to the value the field is currently authoritative for.  That value is
the committed value while an echo is in flight (\texttt{pending}), and the raw
Hub value otherwise.  Honouring the committed value during the window is the
optimistic local echo that closes the durability defect: the display shows what
the user just committed, so a re-focus that lands in the window --- and any edit
begun from it --- builds on the committed value, not on a stale pre-echo one.

There are two honour frames, one idle and one focused, distinguished only by
\texttt{focused}.  They share the same honour expression.  The focused one is
the frame the corrected design adds: it is enabled precisely because focus no
longer begins deferring, so a focused-but-not-yet-edited field still tracks the
Hub.  Deleting it, and making focus defer, is defect~(i).

\begin{schema}{IdleHonour}
  \Delta Field
\where
  focused = zfalse \\
  editing = zfalse \\
  disp' = \IF pending = ztrue \THEN committed \ELSE hub \\
  focused' = focused \\
  editing' = editing \\
  edited' = edited \\
  hub' = hub \\
  pending' = pending \\
  committed' = committed \\
  fires' = fires \\
  lost' = lost \\
  clobbered' = clobbered
\end{schema}

\begin{schema}{FocusedHonour}
  \Delta Field
\where
  focused = ztrue \\
  editing = zfalse \\
  disp' = \IF pending = ztrue \THEN committed \ELSE hub \\
  focused' = focused \\
  editing' = editing \\
  edited' = edited \\
  hub' = hub \\
  pending' = pending \\
  committed' = committed \\
  fires' = fires \\
  lost' = lost \\
  clobbered' = clobbered
\end{schema}

A deferring frame --- \texttt{editing} set --- writes nothing to the display from
the Hub: the buffer is authoritative, so \texttt{disp} changes only by a further
keystroke.  There is no separate schema for it; it is the absence of an honour
frame while \texttt{editing} holds.  That absence is the whole of the no-clobber
guarantee, and its removal --- an echo written to \texttt{disp} regardless of
\texttt{editing} --- is defect~(iii).

One of the frames or the session operations is always enabled: an unfocused
non-deferring field can \texttt{IdleHonour}; a focused non-deferring field can
\texttt{FocusedHonour}; a deferring field is focused and can \texttt{Keystroke}
or \texttt{Commit}; and \texttt{HubEcho} is enabled whenever a commit is in
flight.  A step is therefore always available, which is why the machine cannot
deadlock (Section~\ref{sec:inv}, invariant~5).

% ============================================================
\section{Invariants}
\label{sec:inv}
% ============================================================

Five properties must hold across all reachable states.  They are stated here as
predicates over \texttt{Field}; how each is discharged is given in
Section~\ref{sec:verify}.  None is placed in the \texttt{Field} schema: a
predicate placed there would be enforced as a guard on every successor, silently
disabling any operation that would break it and thereby masking the very defect
this model exists to catch.

\paragraph{Invariant 1 --- commit durability.}
No edit session begins from a stale pre-echo base, so a committed value is never
silently overwritten by a re-focus and edit that started before its echo
arrived.  The bad event --- deferring begun while a commit is in flight and the
base is not the committed value --- latches \texttt{lost}; the invariant is that
it never does:
\[
  lost = zfalse.
\]

\paragraph{Invariant 2 --- honour discipline.}
A field defers only after a genuine edit; it never defers on mere focus.  While
it has not been edited it honours the Hub, so an idle or focused-but-not-edited
field reflects the latest Hub value rather than a frozen snapshot:
\[
  editing = ztrue \implies edited = ztrue.
\]
The display-level reading --- a not-editing field's \texttt{disp} tracks the
value it is authoritative for --- is a corollary: whenever \texttt{editing} is
clear an honour frame is enabled and drives \texttt{disp} to that value, so no
not-editing field is left frozen.  An \texttt{AgentPush} while idle is honoured
on the next \texttt{IdleHonour}; an echo arriving at a focused-not-edited field
is honoured on the next \texttt{FocusedHonour}.

\paragraph{Invariant 3 --- no clobber.}
An echo never overwrites live post-edit text.  A deferring field's display is
moved only by a keystroke, never by a Hub value.  The bad event --- an echo
written to \texttt{disp} while editing --- latches \texttt{clobbered}; the
invariant is that it never does:
\[
  clobbered = zfalse.
\]
This is established by construction: \texttt{HubEcho} frames \texttt{disp}
unchanged, and the only writers of \texttt{disp} are the two honour frames, both
guarded on \texttt{editing = zfalse}, and \texttt{Keystroke}.

\paragraph{Invariant 4 --- commit once.}
A single focus session fires at most once.  With the counter reset on every
focus, and a commit both requiring \texttt{editing} and clearing
\texttt{focused}:
\[
  fires \leq 1.
\]

\paragraph{Invariant 5 --- deadlock freedom.}
No reachable state leaves the machine unable to progress.  The honour frames
partition the not-deferring case by \texttt{focused}; a deferring field can
always \texttt{Keystroke} or \texttt{Commit}; so a step is enabled in every
state.

% ============================================================
\section{Verification}
\label{sec:verify}
% ============================================================

Type-checking is a \texttt{fuzz} pass:
\begin{verbatim}
  fuzz -t docs/input_text_reconciliation.tex
\end{verbatim}

Invariants 1, 2, 3, and 4 are state properties, checked by asking ProB whether
their negation is \emph{reachable}.  A goal that is not found means the invariant
holds on every reachable state:
\begin{verbatim}
  # Invariant 1 -- commit durability (must report: goal NOT found)
  probcli docs/input_text_reconciliation.tex -model_check -nodead \
    -goal "lost=ztrue" -p DEFAULT_SETSIZE 2

  # Invariant 2 -- honour discipline (must report: goal NOT found)
  probcli docs/input_text_reconciliation.tex -model_check -nodead \
    -goal "editing=ztrue & edited=zfalse" -p DEFAULT_SETSIZE 2

  # Invariant 3 -- no clobber (must report: goal NOT found)
  probcli docs/input_text_reconciliation.tex -model_check -nodead \
    -goal "clobbered=ztrue" -p DEFAULT_SETSIZE 2

  # Invariant 4 -- commit once (must report: goal NOT found)
  probcli docs/input_text_reconciliation.tex -model_check -nodead \
    -goal "fires>1" -p DEFAULT_SETSIZE 2
\end{verbatim}

Invariant 5 is the deadlock check over the full reachable state space:
\begin{verbatim}
  # Invariant 5 -- deadlock freedom (must report: no deadlock, no violation)
  probcli docs/input_text_reconciliation.tex -model_check -p DEFAULT_SETSIZE 2
\end{verbatim}

All five return the same verdict at \texttt{DEFAULT\_SETSIZE 3}: the extra value
widens the interleaving of a second edit against an in-flight echo without
adding a new reachable violation.

% ============================================================
\section{Fidelity: reproducing the defects}
\label{sec:fidelity}
% ============================================================

A model that cannot reproduce the known defects proves nothing.  Each defect is
the corrected specification with one part weakened; each then makes the
corresponding invariant's negation reachable.  Below, each variant names the
edit, the goal that changes from \emph{not found} to \emph{found}, and the
shortest offending trace.  Under \texttt{DEFAULT\_SETSIZE 2},
\texttt{VALUE = \{VALUE1, VALUE2\}} with \texttt{v0 = VALUE1}; write $v_0 = v0$
and $v_1$ for the other value.

\paragraph{Defect (i) --- defer on focus (the reported residual).}
This is the current renderer.  Two edits, both away from the corrected design:
\texttt{Focus} sets \texttt{editing' = ztrue} (it snapshots the displayed text
and begins deferring on the mere act of focus, and carries the same durability
latch as \texttt{Keystroke}, now on the \texttt{editing} \texttt{zfalse} $\to$
\texttt{ztrue} transition it introduces); and both honour frames drop the
optimistic echo, writing \texttt{disp' = hub} unconditionally (the renderer's
\texttt{release} dropping the buffer, so an idle frame renders the raw pre-echo
Hub value).  \texttt{FocusedHonour} is thereby dead --- a focused field always
defers --- which is the point.  Goals found: invariant~2
(\texttt{editing=ztrue \& edited=zfalse}) and invariant~1 (\texttt{lost=ztrue}).
Trace for the honour failure:
\begin{enumerate}
  \item \texttt{Init}: \texttt{disp = hub = $v_0$}, not focused, not editing.
  \item \texttt{Focus} in the buggy form: \texttt{focused = ztrue},
    \texttt{editing = ztrue}, \texttt{edited = zfalse}.  This state alone
    violates invariant~2: the field defers although no edit has occurred.
\end{enumerate}
Trace for the durability failure:
\begin{enumerate}
  \item \texttt{Init}; \texttt{Focus}; \texttt{Keystroke($v_1$)}:
    \texttt{disp = $v_1$}, \texttt{editing = ztrue}.
  \item \texttt{Commit}: \texttt{committed = $v_1$}, \texttt{pending = ztrue},
    \texttt{hub = $v_0$} (echo delayed), \texttt{disp = $v_1$},
    \texttt{focused = zfalse}, \texttt{fires = 1}.
  \item \texttt{IdleHonour} in the buggy form: \texttt{disp' = hub = $v_0$} ---
    the display reverts to the stale pre-echo value.
  \item \texttt{Focus} in the buggy form: \texttt{editing} goes \texttt{zfalse}
    $\to$ \texttt{ztrue} while \texttt{pending = ztrue} and \texttt{disp = $v_0$
    $\neq$ $v_1$ = committed} --- deferring begins from a stale base and
    \texttt{lost} latches.  The arriving echo is now deferred for the whole
    session, and the next commit will fire $v_0$, overwriting the committed
    $v_1$.
\end{enumerate}
Under the corrected design, \texttt{Focus} leaves \texttt{editing = zfalse}, and
the honour frames keep \texttt{disp = $v_1$} while \texttt{pending}, so no base
is ever stale.  ProB's own shortest witness for \texttt{lost} takes a different
route to the same stale base: an \texttt{AgentPush} diverges \texttt{hub} from
\texttt{disp}, and a commit fires with no prior keystroke --- itself possible
only because defer-on-focus lets a field be committed without ever being edited
--- after which a raw-hub honour and a re-focus defer from the stale value.  The
loss is thus a property of the defer-on-focus/raw-hub pair, not of one path to
it.

\paragraph{Defect (ii) --- defer after first edit, but honour the raw Hub (the
tempting incomplete fix).}
This is the naive candidate: gate deferring on a real edit --- so invariant~2 is
restored --- but leave the honour frames writing \texttt{disp' = hub}
unconditionally, still dropping the optimistic echo.  It is the fix an informal
argument reaches for, and the model shows it is not enough.  Edit: only the two
honour frames' \texttt{disp'} become \texttt{hub}; \texttt{Focus} and
\texttt{Keystroke} are as corrected.  Goal found: invariant~1
(\texttt{lost=ztrue}); invariant~2 (\texttt{editing=ztrue \& edited=zfalse})
remains \emph{not found}.  Trace:
\begin{enumerate}
  \item \texttt{Init}; \texttt{Focus}; \texttt{Keystroke($v_1$)};
    \texttt{Commit}: \texttt{committed = $v_1$}, \texttt{pending = ztrue},
    \texttt{hub = $v_0$}, \texttt{disp = $v_1$}, \texttt{focused = zfalse}.
  \item \texttt{IdleHonour} in the buggy form: \texttt{disp' = hub = $v_0$} ---
    the field honours the raw pre-echo value and the display goes stale, exactly
    as before, because a not-yet-arrived echo cannot be honoured.
  \item \texttt{Focus}: \texttt{focused = ztrue}, still \texttt{editing =
    zfalse} (correct).
  \item \texttt{Keystroke($v_0$)}: the user types before the echo arrives.
    \texttt{editing} goes \texttt{zfalse} $\to$ \texttt{ztrue} while
    \texttt{pending = ztrue} and \texttt{disp = $v_0$ $\neq$ $v_1$ = committed}
    --- the base is stale and \texttt{lost} latches.  The committed $v_1$ is
    overwritten by an edit begun from $v_0$.
\end{enumerate}
This is the model's central finding.  Deferring after the first edit closes the
\emph{focus-and-sit} route --- a field that sits focused honours the echo when it
arrives --- but it does not close the \emph{type-before-echo} route: a user
faster than the round trip still edits from the stale displayed value.  The
durability defect is not a property of \emph{when} deferring begins but of
\emph{what an idle field displays during the window}.  Only honouring the
committed value locally (the corrected \texttt{IdleHonour}/\texttt{FocusedHonour})
removes the stale base entirely, and the reachability search confirms it:
\texttt{lost=ztrue} is \emph{not found} for the corrected design and
\emph{found}, via the trace above, for this variant.  ProB's shortest witness
reaches it through \texttt{FocusedHonour} rather than \texttt{IdleHonour} --- the
focused-but-not-edited frame itself pulls \texttt{disp} to the stale
\texttt{hub}, then the first keystroke defers from it --- confirming the loss is
independent of whether the stale honour happened while idle or while focused.

\paragraph{Defect (iii) --- echo over live text (the original clobber).}
This is the first defect of the two, the one the commit-on-idle redesign was
meant to prevent: an echo of an earlier edit arriving while a later edit is live.
Edit: \texttt{HubEcho} additionally writes \texttt{disp' = committed} and, when
it does so while \texttt{editing} holds, latches \texttt{clobbered} ---
modelling an echo applied to the buffer without regard to whether the field is
editing.  Goal found: invariant~3 (\texttt{clobbered=ztrue}).  Trace:
\begin{enumerate}
  \item \texttt{Init}; \texttt{Focus}; \texttt{Keystroke($v_1$)};
    \texttt{Commit}: \texttt{committed = $v_1$}, \texttt{pending = ztrue},
    \texttt{hub = $v_0$}.
  \item \texttt{Focus}; \texttt{Keystroke($v_2$)}: a second edit is now live ---
    \texttt{editing = ztrue}, \texttt{disp = $v_2$} --- while the first commit's
    echo is still in flight.
  \item \texttt{HubEcho} in the buggy form: \texttt{disp' = committed = $v_1$}
    while \texttt{editing = ztrue} --- the echo of the first edit overwrites the
    live second edit and \texttt{clobbered} latches.
\end{enumerate}
Under the corrected design, \texttt{HubEcho} frames \texttt{disp} unchanged: a
deferring field's display is the buffer, moved only by a keystroke, so the echo
updates \texttt{hub} silently and the live text stands.

\paragraph{Summary.}
Each defect's weakening makes its invariant's negation reachable; with the
corrected design intact, the four reachability goals (\texttt{lost=ztrue},
\texttt{editing=ztrue \& edited=zfalse}, \texttt{clobbered=ztrue}, and
\texttt{fires>1}) all return \emph{not found} and the deadlock check reports
none.  The difference is the evidence that the model detects the defects the
code must prevent --- and, for defect~(ii), that deferring after the first edit
is \emph{not} the fix on its own: the load-bearing change is honouring the
committed value locally through the echo-latency window.

% ============================================================
\section{From the Model to the Code}
% ============================================================

The corrected design is two changes to the arbiter and its renderer, matching the
two corrected halves of the model.

The first is \emph{defer after the first edit}, not on focus.  The model's
\texttt{Focus} leaves \texttt{editing = zfalse}; only \texttt{Keystroke} sets it.
In the renderer, the buffer is snapshotted and the editing flag set only when
ImGui reports a genuine edit this frame (its \texttt{changed} return, or
\texttt{is\_item\_edited}), not merely when \texttt{is\_item\_active} --- which is
true on mere focus.  The arbiter's \texttt{keep} becomes conditional: it records
the buffer and sets editing when an edit occurred or the field is already
editing, and otherwise leaves the field honouring.  This restores invariant~2.

The second is \emph{honour the committed value through the echo-latency window}.
The model's honour frames render \texttt{committed} while \texttt{pending}, not
the raw Hub value.  In the arbiter, a commit no longer drops the buffer to the
raw Hub value: it records the committed value and the Hub value observed at
commit time, and while not editing \texttt{resolve} returns the committed value
until the Hub value moves off that observed value --- the echo landing, or an
agent override --- at which point it clears the record and returns the Hub value.
This is the model's \texttt{pending}/\texttt{committed} pair and its
\texttt{HubEcho} clearing \texttt{pending}, and it is the change that removes the
stale base entirely, closing invariant~1 including the type-before-echo route
that defer-after-first-edit alone leaves open.

Concretely, in \texttt{InputTextArbiter}
(\texttt{display/renderers/imgui/input\_text\_selection.py}):

\begin{itemize}
  \item Add two slots beside the editing flag, under the element id with fresh
    suffixes (a committed-value slot and a commit-time Hub-value slot), and clear
    them in \texttt{WidgetState.discard\_for} alongside the editing flag so a
    removed field starts clean.  They survive a re-push, like the editing flag,
    because a commit may still be in flight across the resend.
  \item Replace \texttt{keep(text)} with an edit-gated form,
    \texttt{observe(edited, text)}: when \texttt{edited} or the field is already
    editing, set the editing flag and record \texttt{text} as the buffer;
    otherwise do nothing, leaving the field honouring.
  \item Add \texttt{commit(text, hub\_value)}: record \texttt{text} as the
    committed value and \texttt{hub\_value} as the commit-time Hub value.  Do not
    clear the editing flag here; the deactivation path does.
  \item Change \texttt{resolve(hub\_value)}: while editing, return the buffer
    (unchanged); otherwise, if a committed value is recorded and
    \texttt{hub\_value} still equals the recorded commit-time Hub value, return
    the committed value; otherwise clear the recorded slots and return
    \texttt{hub\_value}.
  \item \texttt{release} clears the editing flag and drops the buffer on a
    non-editing frame, but must \emph{not} clear the committed-value slots ---
    they are cleared only when \texttt{resolve} observes the Hub value has moved.
\end{itemize}

In \texttt{InputTextRenderer.render}
(\texttt{display/renderers/input\_text\_renderer.py}), thread the widget's
\texttt{changed} return into the arbiter and record the commit:

\begin{itemize}
  \item Capture the \texttt{changed} flag from \texttt{input\_text\_with\_hint}.
  \item On \texttt{is\_item\_active}, call \texttt{observe(changed, text)} in
    place of \texttt{keep(text)}; otherwise call \texttt{release}.
  \item On \texttt{is\_item\_deactivated\_after\_edit}, both \texttt{fire} the
    \texttt{ValueChanged} and call \texttt{commit(text, elem.value)} so the
    committed value is honoured locally until its echo returns.
\end{itemize}

No other honour- or commit-related behaviour changes.  In particular the
deferring path is untouched: while a real edit is live the buffer remains
authoritative and an echo updates only the Hub value, which is invariant~3
standing as it did.

% ============================================================
\section{Scope and Known Limitations}
\label{sec:scope}
% ============================================================

The reconciliation is by \emph{value equality} alone.  A commit carries no echo
token or version number, so the arbiter recognises an echo only by the Hub value
returning to the value observed at commit time; and one committed / commit-time
Hub-value slot pair records only the \emph{latest} commit.  Two limitations
follow from that single design choice.  Both are transient display artefacts, not
data loss --- the committed value still reaches the Hub and its own echo still
lands --- and both require an event to fall \emph{within one echo round-trip}, so
on localhost, where that window is sub-millisecond, neither is observable in
practice.

\paragraph{Two commits within one round-trip.}
If a field commits twice before the first echo returns, the single slot pair
holds only the second commit and the first commit's Hub value.  While the first
echo is in flight the display can transiently revert to the intermediate
authoritative Hub value --- a flicker between the two committed values --- before
the second commit's own echo lands and the display settles.  Widening the record
to a queue or tagging each commit with a version would remove the flicker at the
cost of the single-slot simplicity; the flicker is preferred while the window is
sub-millisecond.

\paragraph{Agent drive back to the commit-time value.}
An agent that drives the Hub value \emph{back to the exact value observed at
commit time} during the echo-latency window is indistinguishable, under
value-equality reconciliation, from the pending echo of that commit.  The push is
therefore masked --- treated as the awaited echo --- and does not re-render until
the field next moves off that value.  A distinct third value is honoured
immediately (the agent-override case of \texttt{resolve}); only the coincidence
of the \emph{same} value collides.

\paragraph{Relation to the verified model.}
Both limitations live outside this specification's verified scope by
construction.  \texttt{AgentPush} (Section~``The Hub Value'') is guarded to the
settled state --- \texttt{focused = zfalse}, \texttt{editing = zfalse},
\texttt{pending = zfalse} --- so the model never drives the Hub value while a
commit is in flight, and the single \texttt{committed} component admits only one
outstanding commit.  The transitions that exhibit the two artefacts --- a second
commit before an echo, and an agent push \emph{during} \texttt{pending} --- are
thus not reachable in the model, and the five invariants make no claim about
them.  They are documented here, rather than modelled, because closing them would
trade the single-slot / value-equality simplicity for machinery the localhost
round-trip does not justify.

\end{document}
